% ============================================================
% CHAPTER 2: TECHNOLOGY OVERVIEW AND JUSTIFICATION
% ============================================================

\chapter{Technology Overview and Justification}

\section{Introduction to the Technology Ecosystem}

The selection of an appropriate technology stack is a critical architectural decision that fundamentally influences system scalability, development velocity, maintainability, and long-term total cost of ownership (TCO). For this e-commerce platform, the chosen technology ecosystem aligns with modern industry trends while remaining accessible for undergraduate-level implementation and future extensibility.

This chapter provides a comprehensive survey of the core technologies, frameworks, and libraries employed in the system, accompanied by comparative analyses that justify each selection against viable alternatives. The discussion is structured around the classical three-tier architecture: **Presentation Layer** (Frontend), **Application Layer** (Backend/API), and **Data Layer** (Database).

\section{Backend Technology Stack}

\subsection{Runtime Environment: Node.js}

\subsubsection{Technical Overview}

**Node.js** is an asynchronous, event-driven JavaScript runtime built on Chrome's V8 JavaScript engine. First released by Ryan Dahl in 2009, Node.js revolutionized server-side development by enabling JavaScript—traditionally confined to browser environments—to execute on servers \cite{nodejs2024}.

Key architectural characteristics include:

\begin{itemize}
    \item \textbf{Non-Blocking I/O:} Node.js employs a single-threaded event loop with asynchronous callbacks, allowing thousands of concurrent connections without the overhead of multi-threading (context switching, memory consumption).
    
    \item \textbf{V8 Engine:} Google's high-performance JavaScript engine compiles JS directly to native machine code, achieving execution speeds comparable to compiled languages for I/O-bound operations.
    
    \item \textbf{NPM Ecosystem:} The Node Package Manager (npm) hosts over 2 million packages, providing pre-built solutions for authentication, database connectivity, validation, etc., accelerating development.
\end{itemize}

\subsubsection{Justification for Node.js Selection}

A comparative analysis was conducted between Node.js and Python (Django/Flask framework) as candidate backend technologies:

\begin{table}[H]
\centering
\caption{Comparative Analysis: Node.js vs. Python for Backend Development}
\label{tab:nodejs_vs_python}
\begin{tabular}{@{}p{0.25\linewidth}p{0.35\linewidth}p{0.35\linewidth}@{}}
\toprule
\textbf{Criterion} & \textbf{Node.js} & \textbf{Python (Django/Flask)} \\
\midrule
\textbf{Concurrency Model} & Event-driven, non-blocking I/O. Handles 10,000+ concurrent connections efficiently. & Thread-based (Django) or async-capable (Flask with asyncio). Higher memory footprint per connection. \\
\textbf{Performance (I/O)} & Excellent for I/O-bound tasks (database queries, API calls). Benchmarks: ~20,000 req/s for simple APIs. & Moderate. ~5,000-10,000 req/s. Better for CPU-intensive tasks (data science). \\
\textbf{Language Unification} & **Full-stack JavaScript**. Shared codebase (validation schemas, types) between frontend and backend. & Requires context switching between Python (backend) and JavaScript (frontend). \\
\textbf{Ecosystem Maturity} & 2M+ npm packages. Extensive e-commerce libraries (Stripe, PayPal SDKs). & PyPI mature for ML/AI. Less extensive for real-time web features. \\
\textbf{Learning Curve} & Moderate. Asynchronous programming (Promises, async/await) requires understanding. & Low. Synchronous by default, easier for beginners. \\
\textbf{Deployment Simplicity} & Single executable runtime. Docker-friendly. & Requires WSGI server (Gunicorn) + reverse proxy (Nginx). \\
\bottomrule
\end{tabular}
\end{table}

\textbf{Decision Rationale:} Node.js was selected primarily for **full-stack JavaScript unification**, enabling code reuse (e.g., validation schemas shared between client and server) and reducing cognitive load for solo development. Additionally, its event-driven architecture is optimal for e-commerce workloads characterized by high concurrency but low CPU intensity (database reads, API integrations).

\subsubsection{Implementation Evidence}

The project's \texttt{server.js} file initializes the Node.js runtime and configures the Express.js application server:

\begin{lstlisting}[language=JavaScript, caption=Node.js Server Initialization (server.js)]
// ==================================================================
// 1. IMPORTS & CONFIGURATION
// ==================================================================
const express = require('express');
const mongoose = require('mongoose');
const bcrypt = require('bcrypt');
const cors = require('cors');
const jwt = require('jsonwebtoken');
const dotenv = require('dotenv');
const { GoogleGenerativeAI } = require("@google/generative-ai");
const path = require('path');

dotenv.config(); // Load environment variables from .env file

const app = express();
const PORT = process.env.PORT || 3000;

// CORS and JSON Parser Middleware
app.use(cors()); // Enable Cross-Origin Resource Sharing
app.use(express.json()); // Parse JSON payloads

// Static File Serving (Images, Public Assets)
app.use('/images', express.static(path.join(__dirname, 'images')));
app.use('/public', express.static(path.join(__dirname, 'public')));
app.use(express.static(__dirname)); // Serve root directory files

// ==================================================================
// 2. DATABASE CONNECTION
// ==================================================================
const mongoUrl = process.env.MONGO_URL;
if (!mongoUrl) {
    console.error("FATAL: MONGO_URL not configured in .env");
    process.exit(1);
}

mongoose.connect(mongoUrl)
    .then(() => console.log('Database Connected Successfully'))
    .catch((err) => console.error('Database Connection Error:', err));
\end{lstlisting}

\textbf{Technical Observations:}
\begin{itemize}
    \item **Environment Variable Management:** Sensitive credentials (database URL, API keys) are loaded via \texttt{dotenv}, adhering to the **Twelve-Factor App** methodology \cite{twelvefactor}.
    \item **Middleware Configuration:** CORS is enabled globally to allow frontend-backend communication across different origins (essential for deployment scenarios where frontend and backend are hosted separately).
    \item **Static File Serving:** Express's built-in \texttt{static} middleware serves product images and frontend assets, eliminating the need for a separate CDN in the development/demo phase.
\end{itemize}

\subsection{Web Framework: Express.js}

\subsubsection{Technical Overview}

**Express.js** is a minimal and flexible Node.js web application framework providing a robust set of features for building web and mobile applications. It is the de facto standard for Node.js web development, with over 30 million weekly npm downloads \cite{npm2024}.

Core features leveraged in this project:

\begin{itemize}
    \item \textbf{Routing:} Define HTTP endpoints (GET, POST, PUT, DELETE) with pattern matching and parameter extraction.
    \item \textbf{Middleware Pipeline:} Chain reusable functions for authentication, logging, error handling, request parsing.
    \item \textbf{Template Engine Agnostic:} While Express supports templating (EJS, Pug), this project uses it purely as a **JSON API server** with a separate frontend.
\end{itemize}

\subsubsection{RESTful API Design Principles}

The system adheres to **Representational State Transfer (REST)** architectural style, characterized by:

\begin{enumerate}
    \item \textbf{Statelessness:} Each HTTP request contains all necessary information (authentication token, request payload). The server does not maintain client session state, enabling horizontal scalability.
    
    \item \textbf{Resource-Based URLs:} Endpoints represent resources (nouns), not actions (verbs).
    \begin{itemize}
        \item Correct: \texttt{GET /api/products/:id}
        \item Incorrect: \texttt{GET /api/getProductById}
    \end{itemize}
    
    \item \textbf{HTTP Method Semantics:}
    \begin{itemize}
        \item GET: Retrieve resource (idempotent)
        \item POST: Create new resource
        \item PUT: Update existing resource (idempotent)
        \item DELETE: Remove resource (idempotent)
    \end{itemize}
    
    \item \textbf{JSON as Data Interchange Format:} All request/response bodies use JSON, leveraging its ubiquity and native JavaScript support.
\end{enumerate}

\subsubsection{Authentication Middleware Implementation}

A critical security feature is the JWT verification middleware, which protects routes requiring authentication:

\begin{lstlisting}[language=JavaScript, caption=JWT Authentication Middleware (server.js)]
// ==================================================================
// 5. MIDDLEWARES
// ==================================================================

/**
 * Middleware to verify JWT token from request headers.
 * Extracts token from 'token' header, validates signature,
 * and attaches decoded user payload to req.user.
 * 
 * @param {Object} req - Express request object
 * @param {Object} res - Express response object
 * @param {Function} next - Express next middleware function
 */
const verifyToken = (req, res, next) => {
    const authHeader = req.headers.token;
    if (authHeader) {
        const token = authHeader.split(" ")[1]; // Extract "Bearer <TOKEN>"
        jwt.verify(token, process.env.JWT_SECRET, (err, userPayload) => {
            if (err) {
                return res.status(403).json({ 
                    message: "Token is invalid or expired!" 
                });
            }
            req.user = userPayload; // Attach user info to request
            next(); // Proceed to route handler
        });
    } else {
        return res.status(401).json({ 
            message: "Authentication required. No token provided." 
        });
    }
};

/**
 * Middleware to verify admin role.
 * Chains with verifyToken to ensure user is authenticated AND has admin role.
 */
const verifyAdmin = (req, res, next) => {
    verifyToken(req, res, () => {
        if (req.user.role === 'admin') {
            next();
        } else {
            res.status(403).json({ 
                message: "Admin access required!" 
            });
        }
    });
};

module.exports = { verifyToken, verifyAdmin };
\end{lstlisting}

\textbf{Security Analysis:}
\begin{itemize}
    \item **Stateless Authentication:** JWTs eliminate server-side session storage (Redis/Memcached not required), simplifying horizontal scaling.
    \item **Role-Based Access Control (RBAC):** The \texttt{verifyAdmin} middleware demonstrates basic RBAC, restricting certain endpoints (e.g., order management, inventory updates) to admin users.
    \item **Secret Key Management:** The JWT signing secret (\texttt{JWT\_SECRET}) is loaded from environment variables, never hardcoded, mitigating exposure risks.
\end{itemize}

\textbf{Potential Vulnerability:} Token expiration should be short (e.g., 15 minutes) with refresh token rotation for production systems. The current implementation uses 3-day expiration for demo convenience, which is **not production-grade**.

\section{Database Technology: MongoDB}

\subsection{NoSQL vs. SQL Decision Analysis}

The choice between relational (SQL) and non-relational (NoSQL) databases is a foundational architectural decision. A comparative analysis was conducted:

\begin{table}[H]
\centering
\caption{SQL vs. NoSQL Database Comparison for E-Commerce}
\label{tab:sql_vs_nosql}
\begin{tabular}{@{}p{0.25\linewidth}p{0.35\linewidth}p{0.35\linewidth}@{}}
\toprule
\textbf{Criterion} & \textbf{SQL (PostgreSQL)} & \textbf{NoSQL (MongoDB)} \\
\midrule
\textbf{Schema} & Rigid schema. ALTER TABLE for migrations. & Flexible schema. Documents can have varying fields. \\
\textbf{Data Integrity} & ACID guarantees. Foreign keys enforce referential integrity. & Eventual consistency by default. No foreign keys (manual enforcement). \\
\textbf{Query Capability} & Powerful JOINs across tables. Complex aggregations. & Aggregation pipeline. Embedded documents reduce JOINs. \\
\textbf{Scalability} & Vertical scaling primarily. Horizontal (sharding) complex. & Horizontal sharding native. Auto-sharding in MongoDB Atlas. \\
\textbf{Development Speed} & Slower. Schema design upfront required. & Faster prototyping. Schema evolves with code. \\
\textbf{JavaScript Integration} & Requires ORM (Sequelize, TypeORM) for object mapping. & Native BSON (binary JSON). Mongoose ODM feels natural. \\
\bottomrule
\end{tabular}
\end{table}

\textbf{Decision Rationale:} **MongoDB** was selected for the following strategic reasons:

\begin{enumerate}
    \item \textbf{JavaScript Ecosystem Integration:} MongoDB stores data in BSON (binary JSON), aligning perfectly with Node.js's native JSON handling. The Mongoose ODM provides an elegant, schema-based abstraction.
    
    \item \textbf{Rapid Prototyping:} Flexible schemas accelerate feature development in an academic timeline. Product attributes (e.g., color variants, technical specs) can evolve without database migrations.
    
    \item \textbf{Embedded Documents:} E-commerce data models often have aggregation roots (e.g., Order document containing embedded OrderItems array). MongoDB's document model naturally represents these relationships without JOIN overhead.
    
    \item \textbf{Cloud Services:} MongoDB Atlas offers a generous free tier (512MB storage) with automatic backups, eliminating infrastructure management overhead.
\end{enumerate}

\textbf{Trade-Off Acknowledgment:} MongoDB's eventual consistency and lack of foreign key constraints require **application-level data validation**. For example, when creating an order, the system must manually verify product stock availability—a constraint that SQL databases enforce automatically via foreign keys and transactions.

\subsection{Schema Design with Mongoose ODM}

Mongoose provides schema definitions for MongoDB collections, acting as a bridge between unstructured NoSQL and structured relational concepts.

\subsubsection{User Schema Implementation}

\begin{lstlisting}[language=JavaScript, caption=User Schema Definition (server.js)]
// ==================================================================
// 4. MODELS (SCHEMAS)
// ==================================================================

/**
 * User Schema
 * Stores authentication credentials, role, and loyalty program data.
 */
const userSchema = new mongoose.Schema({
    email: { 
        type: String, 
        required: true, 
        unique: true, 
        lowercase: true // Normalize email to lowercase
    },
    password: { 
        type: String, 
        required: true // Hashed with bcrypt before storage
    },
    role: { 
        type: String, 
        default: 'user' // 'user' or 'admin'
    },
    rank: { 
        type: String, 
        enum: ['Silver', 'Gold', 'VIP'], 
        default: 'Silver' // Loyalty program rank
    },
    points: { 
        type: Number, 
        default: 0 // Loyalty points balance
    },
    totalSpending: { 
        type: Number, 
        default: 0 // Cumulative purchase amount (for rank calculation)
    },
    myVouchers: [{
        code: String,
        discountAmount: Number,
        isUsed: { type: Boolean, default: false },
        redeemedAt: { type: Date, default: Date.now }
    }]
}, { timestamps: true }); // Auto-creates createdAt, updatedAt fields

const User = mongoose.model('User', userSchema);
\end{lstlisting}

\textbf{Design Observations:}
\begin{itemize}
    \item **Email Uniqueness:** The \texttt{unique: true} constraint creates a MongoDB unique index, preventing duplicate registrations at the database level.
    \item **Embedded Subdocuments:** The \texttt{myVouchers} array demonstrates MongoDB's ability to embed related data. This design avoids a separate Vouchers collection with foreign keys, trading normalization for query performance.
    \item **Loyalty System Fields:** \texttt{rank}, \texttt{points}, and \texttt{totalSpending} implement a gamified loyalty program—a value-added feature distinguishing this platform from basic e-commerce demos.
\end{itemize}

\subsubsection{Product Schema with Text Indexing}

\begin{lstlisting}[language=JavaScript, caption=Product Schema with Search Optimization (model/Product.js)]
const mongoose = require('mongoose');

const productSchema = new mongoose.Schema({
    name: { type: String, required: true },
    price: { type: Number, required: true },
    short_description: { type: String },
    spec: { type: String }, // Technical specifications
    image_url: { type: String },
    category: { type: String }, // e.g., 'iPhone', 'Mac', 'iPad'
    stock: { type: Number, default: 100 }
}, { timestamps: true });

// TEXT INDEX for full-text search capability
productSchema.index({
    name: 'text',
    short_description: 'text',
    spec: 'text'
});

const Product = mongoose.model('Product', productSchema);
module.exports = Product;
\end{lstlisting}

\textbf{Performance Optimization:}
\begin{itemize}
    \item **Text Indexes:** The compound text index enables efficient full-text search queries (e.g., searching "macbook pro m3" across name, description, and specs). Without this index, MongoDB would perform a collection scan (\texttt{O(n)} complexity), degrading performance with large datasets.
    \item **Stock Field:** Real-time inventory management requires atomic operations. MongoDB's \texttt{findOneAndUpdate} with \texttt{\$inc} operator ensures thread-safe stock decrements during order placement.
\end{itemize}

\subsubsection{Order Schema: Embedded vs. Referenced Design}

\begin{lstlisting}[language=JavaScript, caption=Order Schema Design (server.js)]
const orderSchema = new mongoose.Schema({
    userId: { 
        type: mongoose.Schema.Types.ObjectId, 
        ref: 'User' // Reference to User collection
    },
    recipientName: { type: String, required: true },
    recipientPhone: { type: String, required: true },
    recipientAddress: { type: String, required: true },
    recipientNotes: String,
    paymentMethod: { type: String, required: true },
    
    // EMBEDDED ARRAY: OrderItems stored within Order document
    items: [{ 
        name: String, 
        price: String, // Snapshot of price at order time
        qty: Number, 
        image: String 
    }],
    
    totalAmountString: String, // Formatted string (e.g., "29,990,000 ₫")
    totalAmountNumeric: Number, // Raw number for calculations
    finalAmount: Number, // After discounts
    appliedVoucher: { type: String, default: null },
    status: { 
        type: String, 
        default: 'Pending' // 'Pending', 'Completed', 'Cancelled'
    }
}, { timestamps: true });

const Order = mongoose.model('Order', orderSchema);
\end{lstlisting}

\textbf{Architectural Decision: Embedded Order Items}

The \texttt{items} array is **embedded** within the Order document rather than referenced to a separate OrderItems collection. This design choice reflects a fundamental trade-off in NoSQL modeling:

\textbf{Advantages of Embedding:}
\begin{itemize}
    \item **Query Performance:** Fetching an order retrieves all associated items in a single database read (\texttt{O(1)}), avoiding JOIN operations.
    \item **Price History Preservation:** Storing item prices as strings at order creation ensures historical accuracy. If the product price changes later, past orders remain unaffected—a critical requirement for financial auditing.
\end{itemize}

\textbf{Disadvantages:}
\begin{itemize}
    \item **Data Duplication:** Product names and images are duplicated across orders, violating normalization principles. However, storage cost is negligible compared to query performance gain.
    \item **Document Size Limit:** MongoDB has a 16MB per-document limit. For typical e-commerce orders (5-10 items), this is not a practical concern.
\end{itemize}

\textbf{Alternative Rejected:} A normalized design with separate OrderItems collection and foreign keys would require \texttt{\$lookup} (MongoDB's JOIN equivalent), introducing latency—unacceptable for order history pages displaying hundreds of orders.

\section{Frontend Technology Stack}

\subsection{Core Technologies: HTML5, CSS3, JavaScript ES6+}

The frontend is deliberately built with **vanilla web technologies** (no React, Vue, Angular) to demonstrate mastery of web fundamentals—a pedagogical choice appropriate for undergraduate education.

\subsubsection{HTML5 Semantic Structure}

Modern HTML5 semantic elements (\texttt{<header>}, \texttt{<nav>}, \texttt{<main>}, \texttt{<section>}) are used extensively for:

\begin{itemize}
    \item **SEO Optimization:** Search engines assign higher weight to content within semantic tags.
    \item **Accessibility:** Screen readers utilize semantic structure for improved navigation.
    \item **Code Readability:** Semantic tags document intent better than generic \texttt{<div>} elements.
\end{itemize}

\subsubsection{CSS Framework: Tailwind CSS}

While vanilla CSS3 could suffice, **Tailwind CSS** (a utility-first framework) was adopted for:

\begin{itemize}
    \item **Development Velocity:** Pre-defined utility classes (\texttt{bg-blue-600}, \texttt{rounded-lg}) eliminate writing custom CSS.
    \item **Consistency:** Design tokens (spacing scale, color palette) ensure visual consistency across pages.
    \item **Responsive Design:** Mobile-first breakpoint prefixes (\texttt{md:}, \texttt{lg:}) simplify responsive layouts.
    \item **Production Optimization:** Tailwind's purge feature removes unused CSS, yielding \textless10KB final bundle.
\end{itemize}

Example from \texttt{store.html}:

\begin{lstlisting}[language=HTML, caption=Tailwind CSS Utility Classes for Product Card]
<div class="bg-[#1c1c1e] rounded-2xl p-6 hover:scale-102 
            transition-transform duration-300 shadow-lg 
            border border-white/10">
    <img src="images/iphone16.png" 
         class="w-full h-48 object-contain mb-4">
    <h3 class="text-white font-bold text-lg mb-2">
        iPhone 16 Pro Max
    </h3>
    <p class="text-gray-400 text-sm mb-4">
        A18 Pro chip. Built for Apple Intelligence.
    </p>
    <button class="w-full bg-blue-600 hover:bg-blue-500 
                   text-white font-medium py-3 rounded-xl 
                   transition-colors">
        Add to Cart
    </button>
</div>
\end{lstlisting}

\textbf{Glass morphism Design Pattern:} The system extensively uses background blur (\texttt{backdrop-filter: blur()}) and transparency to achieve Apple-inspired glassmorphism aesthetics, enhancing perceived quality.

\subsection{Advanced Feature: 3D Product Visualization with Three.js}

To differentiate from conventional e-commerce demos, the landing page (\texttt{index.html}) features an interactive 3D iPhone model rendered using **Three.js**, a WebGL abstraction library.

\subsubsection{Three.js Architecture}

\begin{lstlisting}[language=JavaScript, caption=Three.js Scene Initialization (index.html)]
import * as THREE from 'three';
import { GLTFLoader } from 'three/addons/loaders/GLTFLoader.js';

// 1. Scene Setup
const scene = new THREE.Scene();
scene.fog = new THREE.FogExp2(0x000000, 0.01); // Atmospheric fog

// 2. Camera Configuration
const camera = new THREE.PerspectiveCamera(
    45, // Field of view
    window.innerWidth / window.innerHeight, // Aspect ratio
    0.1, // Near clipping plane
    100  // Far clipping plane
);
camera.position.set(0, 0, 15);

// 3. WebGL Renderer
const renderer = new THREE.WebGLRenderer({ 
    canvas: document.querySelector('#webgl-canvas'), 
    antialias: true, 
    alpha: true 
});
renderer.setSize(window.innerWidth, window.innerHeight);
renderer.setPixelRatio(Math.min(window.devicePixelRatio, 2));
renderer.toneMapping = THREE.ACESFilmicToneMapping;
renderer.toneMappingExposure = 2.5;

// 4. Lighting Setup
const ambientLight = new THREE.AmbientLight(0xffffff, 3);
scene.add(ambientLight);

const keyLight = new THREE.SpotLight(0x2997ff, 250);
keyLight.position.set(15, 15, 15);
scene.add(keyLight);

// 5. Load 3D Model (GLTF format)
const loader = new GLTFLoader();
let iphoneModel = null;

loader.load('/images/iphone_16_pro_max.glb', (gltf) => {
    iphoneModel = gltf.scene;
    iphoneModel.scale.set(2.5, 2.5, 2.5);
    iphoneModel.position.set(0, -1, 0);
    
    // Apply PBR materials for photorealism
    iphoneModel.traverse((child) => {
        if (child.isMesh) {
            child.material.metalness = 1.0;
            child.material.roughness = 0.15;
        }
    });
    
    scene.add(iphoneModel);
});

// 6. Animation Loop
function animate() {
    requestAnimationFrame(animate);
    if (iphoneModel) {
        iphoneModel.rotation.y += 0.005; // Slow rotation
    }
    renderer.render(scene, camera);
}
animate();
\end{lstlisting}

\textbf{Technical Achievements:}
\begin{itemize}
    \item **GLTF Format:** The iPhone model is stored in GL Transmission Format (glTF), the "JPEG of 3D"—an industry-standard format optimized for web delivery.
    \item **Physically-Based Rendering (PBR):** Metalness and roughness properties simulate realistic material interaction with light, achieving photorealistic titanium appearance.
    \item **Performance Optimization:** \texttt{requestAnimationFrame} ensures 60 FPS rendering synchronized with display refresh rate.
\end{itemize}

\textbf{User Interaction:} Mouse movement triggers subtle model rotation, creating an interactive "product showcase" experience reminiscent of Apple's official website.

\section{Third-Party Integrations}

\subsection{Google OAuth 2.0 for Social Authentication}

Implementing social login reduces friction in the registration process. The system integrates **Google OAuth 2.0** via the \texttt{passport-google-oauth20} library.

\begin{lstlisting}[language=JavaScript, caption=Google OAuth Strategy Configuration (server.js)]
const passport = require('passport');
const GoogleStrategy = require('passport-google-oauth20').Strategy;
const session = require('express-session');

// Session configuration for production (HTTPS required)
app.use(session({
    secret: 'apple_store_secret_key',
    resave: false,
    saveUninitialized: true,
    cookie: {
        secure: true, // HTTPS only
        sameSite: 'none', // Cross-site cookie
        maxAge: 24 * 60 * 60 * 1000 // 1 day
    }
}));

app.use(passport.initialize());
app.use(passport.session());

// Google OAuth Strategy
passport.use(new GoogleStrategy({
    clientID: process.env.GOOGLE_CLIENT_ID,
    clientSecret: process.env.GOOGLE_CLIENT_SECRET,
    callbackURL: "https://project3-icy1.onrender.com/auth/google/callback"
}, async (accessToken, refreshToken, profile, done) => {
    try {
        // Find existing user by email
        let user = await User.findOne({ email: profile.emails[0].value });
        
        if (!user) {
            // Create new user if not exists
            const randomPassword = await bcrypt.hash(
                Math.random().toString(36), 10
            );
            user = new User({
                email: profile.emails[0].value,
                password: randomPassword, // Secure random password
                role: 'user',
                rank: 'Silver',
                points: 0
            });
            await user.save();
        }
        return done(null, user);
    } catch (err) {
        return done(err, null);
    }
}));

// OAuth Routes
app.get('/auth/google', 
    passport.authenticate('google', { scope: ['profile', 'email'] })
);

app.get('/auth/google/callback',
    passport.authenticate('google', { failureRedirect: '/login' }),
    (req, res) => {
        // Generate JWT after successful Google auth
        const accessToken = jwt.sign(
            { id: req.user._id, role: req.user.role },
            process.env.JWT_SECRET,
            { expiresIn: "3d" }
        );
        res.redirect(`/?token=${accessToken}`);
    }
);
\end{lstlisting}

\textbf{Security Considerations:}
\begin{itemize}
    \item **OAuth 2.0 Authorization Code Flow:** The system uses the most secure OAuth flow, where the client secret is never exposed to the browser.
    \item **Automatic Account Linking:** If a user previously registered with email/password and later uses Google OAuth with the same email, the system links to the existing account, preventing duplicates.
\end{itemize}

\subsection{AI Chatbot: Google Gemini API Integration}

Modern e-commerce platforms increasingly employ AI assistants for customer support. This project integrates **Google Gemini 2.5 Flash**, a state-of-the-art large language model (LLM) optimized for low latency.

\begin{lstlisting}[language=JavaScript, caption=AI Chatbot Implementation (server.js)]
const { GoogleGenerativeAI } = require("@google/generative-ai");

const apiKey = process.env.GEMINI_API_KEY;
const genAI = new GoogleGenerativeAI(apiKey);
const model = genAI.getGenerativeModel({ model: "gemini-2.5-flash" });

app.post('/api/chat', async (req, res) => {
    const userMessage = req.body.message;
    
    // Optional user authentication
    let userId = null;
    const authHeader = req.headers.token;
    if (authHeader) {
        try {
            const token = authHeader.split(" ")[1];
            userId = jwt.verify(token, process.env.JWT_SECRET).id;
        } catch (e) { /* Guest user */ }
    }
    
    // Build context from user data
    let contextData = { recent_orders: [], available_products: [] };
    
    if (userId) {
        const orders = await Order.find({ userId })
                                .sort({ createdAt: -1 }).limit(5);
        contextData.recent_orders = orders.map(o => ({
            id: o._id.toString().slice(-6).toUpperCase(),
            status: o.status,
            total: o.finalAmount.toLocaleString('vi-VN') + 'đ'
        }));
    }
    
    const products = await Product.find({ stock: { $gt: 0 } })
                                  .select('name price category').limit(50);
    contextData.available_products = products.map(p => ({
        name: p.name,
        price: p.price.toLocaleString('vi-VN') + 'đ',
        category: p.category
    }));
    
    // Generate AI response with context
    const systemPrompt = `
        YOU ARE: AI assistant for Apple Store e-commerce.
        DATA:
        - User Orders: ${JSON.stringify(contextData.recent_orders)}
        - Available Products: ${JSON.stringify(contextData.available_products)}
        TASK: Answer concisely about orders, products, prices in Vietnamese.
        User question: "${userMessage}"
    `;
    
    try {
        const result = await model.generateContent(systemPrompt);
        const response = await result.response;
        const replyText = response.text();
        
        res.json({ type: 'text', reply: replyText });
    } catch (error) {
        res.status(500).json({ 
            reply: "System error. Contact: 0962923329" 
        });
    }
});
\end{lstlisting}

\textbf{Implementation Highlights:}
\begin{itemize}
    \item **Context-Aware Responses:** The AI has access to real-time product catalog and user order history, enabling personalized responses (e.g., "Your last order #ABC123 is in Shipped status").
    \item **Fallback Mechanism:** If Gemini API fails (quota exceeded, network error), the system gracefully returns an error message with customer service contact, maintaining UX quality.
    \item **Performance Optimization:** Simple price queries are handled via regex pattern matching without API calls, reducing latency and costs.
\end{itemize}

\section{Supporting Tools and Libraries}

\begin{table}[H]
\centering
\caption{Complete Technology Stack Summary}
\label{tab:tech_stack_summary}
\small
\begin{tabular}{@{}p{0.20\linewidth}p{0.30\linewidth}p{0.40\linewidth}@{}}
\toprule
\textbf{Layer} & \textbf{Technology/Library} & \textbf{Purpose and Justification} \\
\midrule
\textbf{Backend Runtime} & Node.js v18.16.0 LTS & Event-driven, non-blocking I/O for scalable APIs \\
\textbf{Web Framework} & Express.js v5.1.0 & Minimalist, middleware-based routing \\
\textbf{Database} & MongoDB v8.15.1 (Mongoose ODM) & Flexible schema, BSON/JSON alignment \\
\textbf{Authentication} & jsonwebtoken v9.0.2 & Stateless JWT token generation/verification \\
 & passport-google-oauth20 v2.0.0 & Google OAuth 2.0 social login \\
\textbf{Security} & bcrypt v5.1.1 & Password hashing (10 salt rounds) \\
 & cors v2.8.5 & Cross-Origin Resource Sharing config \\
\textbf{AI Integration} & @google/generative-ai v0.24.1 & Gemini 2.5 Flash LLM API client \\
\textbf{Frontend Base} & HTML5, CSS3, JavaScript ES6+ & Semantic structure, modern syntax \\
\textbf{CSS Framework} & Tailwind CSS (CDN) & Utility-first styling, rapid prototyping \\
\textbf{3D Graphics} & Three.js v0.176.0 & WebGL-based 3D model rendering \\
\textbf{Session Management} & express-session v1.18.0 & OAuth callback session handling \\
\textbf{Environment Variables} & dotenv v17.2.2 & Secure credential management \\
\textbf{Development} & nodemon v3.1.10 & Auto-restart on code changes \\
 & concurrently v9.1.2 & Run frontend + backend simultaneously \\
\bottomrule
\end{tabular}
\end{table}

\section{Architecture Pattern: Modular Monolith}

The system is architected as a **Modular Monolith**—a hybrid approach balancing the simplicity of monolithic deployment with the modularity of microservices.

\subsection{Monolithic vs. Microservices vs. Modular Monolith}

\begin{table}[H]
\centering
\caption{Architecture Pattern Comparison}
\label{tab:architecture_patterns}
\begin{tabular}{@{}p{0.25\linewidth}p{0.22\linewidth}p{0.22\linewidth}p{0.22\linewidth}@{}}
\toprule
\textbf{Criterion} & \textbf{Monolith} & \textbf{Microservices} & \textbf{Modular Monolith} \\
\midrule
\textbf{Deployment} & Single process & Multiple services & Single process \\
\textbf{Scalability} & Vertical only & Horizontal per service & Horizontal (entire app) \\
\textbf{Development Complexity} & Low & High (distributed systems) & Moderate \\
\textbf{Latency} & Low (in-process calls) & Higher (network calls) & Low \\
\textbf{Technology Diversity} & Single stack & Polyglot possible & Single stack \\
\textbf{Suitable For} & MVPs, small teams & Large orgs, high scale & SMEs, growing startups \\
\bottomrule
\end{tabular}
\end{table}

\textbf{Selected Approach:} **Modular Monolith** was chosen to maintain **logical separation** (authentication module, product module, order module) within a **single codebase and deployment unit**. This provides:

\begin{itemize}
    \item **Simplified DevOps:** One Docker container, one database connection pool.
    \item **Future Microservices Migration Path:** Well-defined module boundaries facilitate future extraction into separate services if traffic scales beyond monolith capabilities.
    \item **Academic Feasibility:** Reduces operational complexity, allowing focus on application logic rather than distributed systems orchestration (Kubernetes, service mesh).
\end{itemize}

\section{Summary}

This chapter presented a comprehensive analysis of the technology stack underpinning the e-commerce platform. Key decisions—Node.js for backend, MongoDB for database, vanilla frontend with Three.js enhancements—were justified through comparative analyses and aligned with modern industry practices.

The next chapter (System Analysis) will translate these technological foundations into concrete functional and non-functional requirements, use cases, and system constraints.

% End of Chapter 2
